Under the global group $S(U(2) \times U(3))$ of the minimal massive composite: unconditionally, the pair-level generators transform as $(2,3)$ with $Y = 1/6$, the quark doublet, and the singlets of a single conformal spinor carry the hypercharges $0, -\tfrac23, \tfrac13, 1$, the latter as a computation on the Fock space of a single conformal spinor, which is not a subspace of the tensor state space. Under Hypothesis (S) (\claimref{standard-model.hyp-s-slot-statistics}) (the slots of the conformal spinors are fermionic modes on the same footing as the pairs, so the composite's state space is the five-mode Fock space), and only under it, the lepton doublet $(2,1)$ with $Y = -\tfrac12$ exists, and, given the content (the two doublets and the singlets listed in one handedness, with a fourth singlet of hypercharge zero or none) and the normalisation $Y_{Q} = 1/6$, the hypercharge spectrum of the generation is unique from the four anomaly conditions, the pair $\{Y_{u^{c}}, Y_{d^{c}}\}$ as an unordered set; the fourth-singlet and common-handedness hypotheses are both necessary there, and without the lepton doublet the anomaly conditions are different equations and the uniqueness is not claimed. Chirality is the parity grading of the Fock module, whose two halves carry inequivalent content. Hypothesis (S) is a structural premise, not a theorem of the series, and every statement of this paper about the charged leptons carries it. \emph{Provenance:} \cite[\S4-\S5, in particular Prop.~5.3, Thm.~5.5, Prop.~5.6]{ButlerT12} (\claimref{standard-model.chirality-two-halves-spinor-module-conjugation}, \claimref{standard-model.quarks-leptons-two-level-fermion-content}, \claimref{standard-model.hypercharge-global-form-pair-counting-consistency}).
